\section{不平衡样本集问题}

\begin{definition}
	在分类问题中，若观测数据集中不同类别的样本数量相差悬殊，则称该数据集存在\gls{Class-imbalancedDataset}问题。此时，样本数量较多的类别称为\gls{MajorityClass}，样本数量较少的类别称为\gls{MinorityClass}。
\end{definition}
不平衡样本集问题的根源在于，模型训练通常以最小化经验风险为目标，而经验风险是各观测样本损失的平均。当不同类别的样本数量相差悬殊时，多数类样本对经验风险的贡献会占据主导地位，因此，即使模型在少数类上的预测表现很差，只要能够较好地拟合多数类，仍可能获得较小的总体经验风险。换言之，经验风险最小化本身并不会保证模型对各类别具有相近的识别能力，而可能使训练过程系统性地偏向多数类。\par
例如，在一个二分类任务中，观测数据集包含$100$个样本，其中$99$个属于多数类$0$，仅有$1$个属于少数类$1$。若模型对任意样本的输出都为$0$，则在$0$-$1$损失下（预测正确损失函数值为$0$，预测错误损失函数值为$1$），其经验风险仅为$0.01$。因此，从总体经验风险的角度看，该模型已经取得了看似良好的预测结果。然而，它完全无法识别任何少数类样本。\par
类别不平衡现象广泛存在于实际分类任务中，并且在许多应用场景下，样本数量较少的类别恰恰对应我们真正关心的事件。例如，疾病患者、欺诈交易和设备故障通常只占全部观测中的很小比例，但识别这些少数类样本往往正是建模任务的核心目标。因此，一个仅通过良好预测多数类而获得较低经验风险的模型，即使在整体意义上表现良好，也可能完全失去实际应用价值，因为它没有学习到任务中最重要的少数类模式。\par
在接下来的讨论中，我们将主要涉及二分类任务，并且默认将少数类记为正类。

\subsection{评价指标}
类别不平衡不仅会影响模型的训练过程，也会影响模型性能的评价。当不同类别的样本数量相差悬殊时，一些在一般分类任务中常用的评价指标可能无法准确反映模型对少数类的识别能力。因此，在讨论不平衡样本集的评价问题时，一个自然的思路是重新考察此前介绍的各类评价指标，分析它们在类别不平衡情形下是否仍然具有合理的解释。对于容易受到类别比例影响的指标，则需要考虑相应的修正形式或采用新的评价指标。\par
在类别不平衡的数据集中，准确率显然容易受到多数类样本的主导。为避免不同类别的样本数量直接决定评价结果，可以分别计算模型在各类别上的识别准确率，再对其进行等权平均，从而减弱类别数量不平衡对评价结果的影响。。
\begin{definition}
	对于二分类任务，称：
	\begin{equation*}
		\operatorname{BalancedAccuracy}=\frac{1}{2}\left(\frac{\operatorname{TP}}{\operatorname{TP}+\operatorname{FN}}+\frac{\operatorname{TN}}{\operatorname{TN}+\operatorname{FP}}\right)
	\end{equation*}
	为模型的\gls{BalancedAccuracy}。
\end{definition}
由于默认将少数类视为正类，因此精确率、召回率、$\operatorname{F1}$分数、$\operatorname{F}_{\beta}$分数以及$\operatorname{AUPRC}$都能够直接反映模型在少数类上的预测表现。精确率刻画模型预测为少数类的样本中有多少真正属于少数类，召回率刻画实际少数类样本中有多少被模型成功识别，$\operatorname{F1}$分数和$\operatorname{F}_{\beta}$分数则综合衡量精确率与召回率，而$\operatorname{AUPRC}$进一步描述不同分类阈值下精确率与召回率之间的整体权衡。这些指标更加关注正类的识别效果，因此通常适用于类别不平衡问题。

\subsection{代价敏感学习}
对于不平衡分类问题，一种直接的处理思路是重新调整不同类别样本在损失函数中的权重，以抵消类别样本数量差异对经验风险的影响，使各类别的预测误差能够在训练目标中得到更为均衡的体现。这类通过赋予不同类型样本的预测错误以不同代价，从而改变模型优化目标的方法称为\gls{CostSensitiveLearning}。\par
具体而言，设分类问题共有$K$个类别，第$k$类样本数量为$n_k$，为该类别赋予权重$w_k>0$，则代价敏感学习一般的经验风险被定义为：
\begin{equation*}
	\hat{R}_w(\theta)\coloneq\sum_{i=1}^{n}w_{y_i}L[y_i,f_\theta(x_i)]
\end{equation*}\par
一种常用的权重设定方式是令类别权重与该类别的样本数量成反比，即：
\begin{equation*}
	w_k\propto\frac{1}{n_k}
\end{equation*}
此时第$k$类样本在加权经验风险中的总权重满足：
\begin{equation*}
	n_kw_k\propto1
\end{equation*}
因此，不同类别虽然具有不同的样本数量和单样本权重，但其在总体损失中具有相同的总权重，从而消除了类别频率本身对经验风险贡献所造成的失衡。

\subsection{样本平衡类方法}
与代价敏感学习直接调整不同类别样本在损失函数中的贡献不同，处理不平衡样本集的另一类基本思路是从训练数据本身出发，通过改变不同类别样本在训练过程中的出现频率，使模型所观察到的经验类别分布更加均衡。这类方法统称为\textbf{样本平衡类方法}。其核心思想是，在不改变原始学习任务和损失函数形式的前提下，通过对多数类样本进行删减、对少数类样本进行重复或合成，重新构造用于模型训练的数据集，从而减弱多数类样本因数量占优而对经验风险及参数更新产生的主导作用。
\subsubsection{欠采样与过采样}
设分类问题共有$K$个类别，第$k$类样本数量为$n_k$。样本平衡最直接的做法是通过重新采样改变不同类别样本参与模型训练的频率，从而调整训练数据中的经验类别分布。根据调整方式的不同，可以分为\gls{UnderSampling}与\gls{OverSampling}。\par
欠采样通过减少多数类样本的数量来缓解类别之间的样本规模差异。最简单的随机欠采样方法从多数类样本中随机选取部分观测用于训练。对于二分类问题，设多数类和少数类的样本数量分别为$n_{\mathrm{maj}}$和$n_{\mathrm{min}}$，可以从多数类中抽取$n_{\mathrm{maj}}'<n_{\mathrm{maj}}$个样本，使$n_{\mathrm{maj}}'$与$n_{\mathrm{min}}$之间的差异缩小。\par
欠采样能够减少多数类对训练目标的主导作用，同时直接降低训练数据规模，因此在多数类样本数量极大时还具有降低存储与计算成本的优势。其主要缺点在于部分已经观测到的样本会被舍弃。如果多数类具有复杂的内部结构，或者被删除的样本包含重要的决策边界信息，则过强的欠采样可能造成有效信息损失，并降低模型的泛化能力。\par
与之相对，过采样不删除原有观测，而是通过提高少数类样本在训练过程中的出现频率来改善类别不平衡。最简单的随机过采样方法从少数类样本中进行有放回抽样，使其有效样本数量增加。例如，可以将少数类从$n_{\mathrm{min}}$个样本重复采样至$n_{\mathrm{min}}'$个，使其与多数类的样本数量更加接近。\par
随机过采样避免了欠采样所带来的信息丢失，但简单复制已有样本并不会产生新的观测信息。少数类中的同一观测可能被模型反复使用，从而增加对有限少数类样本过拟合的风险；与此同时，重复采样还会增大训练数据规模和计算开销。
\subsubsection{SMOTE}
随机过采样通过重复已有少数类样本提高其在训练过程中的出现频率，但重复观测本身并不会引入新的特征信息。为缓解这一问题，\gls{SMOTE}不再直接复制少数类样本，而是利用少数类样本在特征空间中的局部邻域结构生成新的合成样本。\par
设$x_i$为某个少数类样本，并在少数类样本中寻找其$k$个近邻。从这些近邻中选取一个样本$x_j$，再生成随机变量：
\begin{equation*}
	\lambda\sim\operatorname{Uniform}(0,1)
\end{equation*}
则SMOTE生成的新样本为：
\begin{equation*}
	x_{\mathrm{new}}=x_i+\lambda(x_j-x_i)
\end{equation*}
因此，$x_{\mathrm{new}}$位于$x_i$与$x_j$之间的线段上，并被赋予与$x_i$相同的类别标签。通过对不同的少数类样本及其近邻重复这一过程，可以逐步扩充少数类训练样本。\par
SMOTE隐含的基本假设是：相邻少数类样本之间的区域仍然属于少数类，因此可以通过局部线性插值生成新的少数类样本。相较于随机过采样直接复制已有观测，这种方式能够扩展少数类样本在特征空间中的覆盖范围。这一假设也是SMOTE的主要局限。当少数类靠近决策边界或少数类本身具有复杂的非线性结构时，两个少数类样本之间的插值区域未必仍然属于少数类，从而可能生成不合理的合成样本并干扰决策边界。

\subsection{Balanced Bagging}
Bagging通过对训练数据进行重复采样并训练多个基学习器，再将各基学习器的预测结果进行集成，利用基学习器之间的异质性去提高模型的预测精度（\cref{prop:TreeEnsemble}(2)）。对于不平衡分类问题，可以进一步将类别平衡操作嵌入Bagging的采样过程中，使每个基学习器均在相对平衡的数据集上训练。这类方法通常称为\textbf{Balanced Bagging}。\par
设二分类训练集为：
\begin{equation*}
	\mathcal{D}=\mathcal{D}_{+}\cup\mathcal{D}_{-}
\end{equation*}
其中$n_{+}\ll n_{-}$。对于第$b$个基学习器，Balanced Bagging保留少数类样本，并从多数类中随机抽取与少数类数量相近的样本，构造平衡训练集$\mathcal{D}^{(b)}$并训练$f^{(b)}$，其最终模型为$f^{(1)},f^{(2)},\dots,f^{(B)}$经过多数投票得到的结果。\par
这种方法通常采用多数类欠采样，而不是少数类过采样。原因在于单次欠采样虽然会丢失多数类信息，但Bagging中的不同基学习器可以使用不同的多数类子集，从而在整个集成模型中逐步利用更多多数类观测，因此Bagging恰好能够缓解欠采样最主要的信息损失问题。相比之下，随机过采样主要是在不同训练集中重复有限的少数类样本，并不会提供新的观测信息，还会增大每个基学习器的训练规模，并可能使不同基学习器反复依赖相同的少数类样本，提高了基学习器之间的相关性，从而无法达到理想的预测精度。

\subsection{困难负样本挖掘}
\begin{definition}
	设二分类训练集为$\mathcal{D}=\mathcal{D}_{+}\cup\mathcal{D}_{-}$，其中$\mathcal{D}_{+}$与$\mathcal{D}_{-}$分别为正类和负类样本集。对于给定模型$f_{\theta}$，定义负类样本$(x_i,0)\in\mathcal{D}_{-}$的分类困难程度为：
	\begin{equation*}
		d_i(\theta)\coloneq L[0,f_{\theta}(x_i)]
	\end{equation*}
	给定阈值$\tau$，若$d_i(\theta)>\tau$，则称$(x_i,0)$为模型$f_{\theta}$的困难负样本。从多数类中选择困难负样本并用于后续训练的方法称为\gls{HardNegativeMining}。
\end{definition}
困难负样本挖掘的核心思想是：多数类样本虽然数量庞大，但并非所有样本对模型训练都具有同等价值。大量能够被当前模型轻易正确分类的负类样本对决策边界的进一步修正作用有限，而那些容易被误判、损失较大的负类样本更能反映模型当前尚未解决的分类困难。因此，与随机欠采样等概率地丢弃多数类样本不同，困难负样本挖掘利用当前模型主动筛选更有训练价值的负类样本，使模型将有限的训练能力集中于难以区分的区域。\par
由于样本的困难程度由当前模型决定，困难负样本并不是固定不变的。随着模型参数更新，原有困难负样本的损失可能降低，而其他负类样本可能成为新的困难样本。因此，困难负样本挖掘通常与模型训练交替进行，即利用当前模型筛选困难负样本，再使用这些样本更新模型，并重复这一过程，直至满足停止准则。
\subsubsection{机器学习中的困难负样本挖掘}
\begin{method}
	设第$t$轮训练所使用的负类样本集为$\mathcal{S}_{-}^{(t)}\subseteq\mathcal{D}_{-}$。对于非神经网络的传统机器学习模型，困难负样本挖掘的一般过程为：
	\begin{enumerate}
		\item 从$\mathcal{D}_{-}$中选取部分负类样本，构造初始负类训练集$\mathcal{S}_{-}^{(0)}$；
		\item 对$t=0,1,2,\ldots$，使用$\mathcal{D}_{+}\cup\mathcal{S}_{-}^{(t)}$训练模型$f_{\theta^{(t)}}$；
		\item 使用$f_{\theta^{(t)}}$评估负类样本的困难程度$d_i(\theta^{(t)})$，并据此确定当前模型下需要重点用于后续训练的困难负样本；
		\item 根据当前负类训练集以及新识别出的困难负样本构造下一轮负类训练集$\mathcal{S}_{-}^{(t+1)}$，随后进入下一轮训练。
	\end{enumerate}
	当验证集上的评价指标不再明显改善，或达到预设的最大迭代轮数时停止。
\end{method}
设$n_{+}=|\mathcal{D}_{+}|$且$s_t=|\mathcal{S}_{-}^{(t)}|$。若还希望正类与选中负类对训练目标具有相同的总贡献，可以在第$t$轮使用如下经验风险：
\begin{equation*}
	\hat{R}_{\mathrm{HNM}}^{(t)}(\theta)\coloneq\frac{1}{n_{+}}\sum_{(x_i,1)\in\mathcal{D}_{+}}L[1,f_{\theta}(x_i)]
	+\frac{1}{s_t}\sum_{(x_i,0)\in\mathcal{S}_{-}^{(t)}}L[0,f_{\theta}(x_i)]
\end{equation*}\par
负类训练集$\mathcal{S}_{-}^{(t)}$的更新方式并不唯一。最直接的做法是在已有负类训练集上不断加入新识别出的困难负样本，使训练集随迭代逐渐扩增；也可以在每一轮根据当前模型重新选择负类样本，从而完全重构$\mathcal{S}_{-}^{(t+1)}$。此外，还可以采用介于两者之间的策略，例如保留部分已有困难负样本，同时移除已经能够被模型轻易识别的样本并补充新的困难负样本。具体采用何种更新方式取决于具体场景。
\subsubsection{深度学习中的困难负样本挖掘}
设模型给出正类概率$q_{\theta}(x_i)\in(0,1)$，记：
\begin{equation*}
	p_{ti}=
	\begin{cases}
		q_{\theta}(x_i), &y_i=1 \\
		1-q_{\theta}(x_i), &y_i=0
	\end{cases}
	\quad
	\alpha_{ti}=
	\begin{cases}
		\alpha,& y_i=1, \\
		1-\alpha,& y_i=0
	\end{cases}
\end{equation*}
其中$\alpha\in(0,1)$为类别权重参数。
\begin{definition}
	设聚焦参数$\gamma\geqslant0$。称：
	\begin{equation*}
		L_{\mathrm{FL}}[y_i,q_{\theta}(x_i)]
		\coloneq-\alpha_{ti}(1-p_{ti})^{\gamma}\ln p_{ti}
	\end{equation*}
	为样本$(x_i,y_i)$的\gls{FocalLoss}。
\end{definition}
Focal loss在交叉熵前加入调节因子$(1-p_{ti})^{\gamma}$。当样本容易分类时$p_{ti}$接近$1$，其损失被显著抑制；当样本难以分类时$p_{ti}$较小，其损失得到保留。因此，Focal loss无需显式筛除容易样本，即可使训练更加关注困难样本。
\subsubsection{困难负样本挖掘的缺陷}
困难负样本挖掘依赖当前模型对样本困难程度的判断，因此其效果容易受到初始模型和迭代过程的影响。如果模型早期的决策边界存在较大偏差，所选出的困难负样本也可能缺乏代表性，并进一步影响后续训练。此外，高损失样本并不一定是真正有信息的困难样本。错误标签、异常观测以及类别重叠区域中的样本同样可能产生较大损失，从而被反复选中并获得过高的训练关注。与此同时，过度强调困难负样本会削弱模型对多数类整体分布的利用，若负样本选择过于集中，还可能降低模型对其他多数类区域的泛化能力。

